\documentclass[11pt]{amsart}

\usepackage[T1]{fontenc}
\usepackage{amsmath,amssymb,amsthm,mathtools}
\usepackage[a4paper,left=25mm,right=25mm,top=23mm,bottom=25mm]{geometry}
\usepackage[hidelinks]{hyperref}
\usepackage{microtype}
\microtypesetup{expansion=false}

\newtheorem{theorem}{Theorem}[section]
\newtheorem{proposition}[theorem]{Proposition}
\newtheorem{lemma}[theorem]{Lemma}
\newtheorem{corollary}[theorem]{Corollary}
\theoremstyle{remark}
\newtheorem{remark}[theorem]{Remark}

\newcommand{\Z}{\mathbb Z}
\newcommand{\Q}{\mathbb Q}
\newcommand{\Sym}{\operatorname{Sym}}
\newcommand{\Hom}{\operatorname{Hom}}
\newcommand{\Ext}{\operatorname{Ext}}
\newcommand{\im}{\operatorname{im}}
\newcommand{\gm}{\gamma}
\newcommand{\dlt}{\delta}
\newcommand{\aug}{\varpi}

\title[The fifth dimension subring]{The fifth dimension subring:\
the class-three base case}
\author{}
\date{August 2026}

\begin{document}

\begin{abstract}
We prove the reduced statement underlying
$\delta _5(L)\subseteq\gamma _4(L)$.  Let $L$ be a finite Lie ring
whose additive group is a $2$-group, assume $\gamma _4(L)=0$, and assume
$C=\gamma _3(L)$ is cyclic.  We prove $\delta _5(L)=0$.  The proof is
separated into two independent assertions.  The normal-form theorem
extracts from a PBW witness a class in the first derived symmetric
square of the additive extension
$0\to\gamma _2(L)/C\to L/C\to L/\gamma _2(L)\to0$.  The vanishing
theorem says that the resulting character vanishes on every class
coming from the middle term.  We include the complete class-three PBW
ledger and all steps of the functorial vanishing argument.
\end{abstract}

\maketitle

\section{Statement and notation}

For a Lie ring $L$, write $U(L)$ for its universal enveloping algebra,
$\aug(L)$ for the augmentation ideal of $U(L)$, and
\[
 \dlt_r(L)=L\cap\aug(L)^r,
 \qquad
 \gm_1(L)=L,
 \qquad
 \gm_{r+1}(L)=[L,\gm_r(L)].
\]
The intersection defining $\dlt_r(L)$ is taken through the canonical
map $L\to U(L)$.

We use two published facts from~\cite[Theorem~4.5 and
Proposition~4.2]{BP}.  First,
$\dlt_3(L)=\gm_3(L)$ for every Lie ring.  Second, if
$F\twoheadrightarrow L=F/\mathcal R$ is a free presentation,
$\mathfrak r=U(F)\mathcal R U(F)$, and $f\in F$ maps to $a\in L$, then
\begin{equation}\label{eq:free-criterion}
 a\in\dlt_r(L)
 \quad\Longleftrightarrow\quad
 f\in \aug(F)^r+\mathfrak r.
\end{equation}

We first prove the precise reduced statement.

\begin{theorem}[reduced base case]\label{thm:reduced}
Let $L$ be a finite Lie ring whose additive group is a $2$-group.
Assume
\[
                 \gm_4(L)=0,
 \qquad
 C:=\gm_3(L)=\langle z\rangle\cong\Z/q,
 \qquad q=2^s.
\]
Then $\dlt_5(L)=0$.
\end{theorem}

The hypothesis $L''=0$ need not be added: $L''=[\gm_2(L),\gm_2(L)]
\subseteq\gm_4(L)=0$.  Put
\[
 H=L/C,
 \qquad V=L/\gm_2(L),
 \qquad W=\gm_2(L)/C.
\]
The underlying additive groups form an exact sequence
\begin{equation}\label{eq:extension}
             0\longrightarrow W\longrightarrow H
             \xrightarrow{p}V\longrightarrow0.
\end{equation}
Since $[L,\gm_2(L)]=C$ and $[C,L]=0$, the bracket induces a bilinear
map
\begin{equation}\label{eq:g}
 g:V\otimes W\longrightarrow C,
 \qquad g(\bar x,\bar w)=[x,w].
\end{equation}
It is independent of both choices of lifts: changing $x$ by an
element of $\gm_2(L)$ changes the bracket by an element of
$[\gm_2(L),\gm_2(L)]\subseteq\gm_4(L)$, while changing $w$ by an
element of $C$ changes it by an element of $[L,C]=\gm_4(L)$.

We identify $C$ with the subgroup $\frac1q\Z/\Z$ of $\Q/\Z$ by
$z\mapsto1/q+\Z$.  The proof of Theorem~\ref{thm:reduced} will be the
composition of the following two statements.

\begin{theorem}[normal form, informal version]\label{thm:nf-informal}
There is a character
$\eta:L_1\Sym^2(V)\to\Q/\Z$ such that, whenever
$\zeta z\in\dlt_5(L)$, there is
$\xi\in L_1\Sym^2(H)$ satisfying
\[
 \frac{\zeta}{q}+\Z
   =\eta\bigl(L_1\Sym^2(p)(\xi)\bigr).
\]
\end{theorem}

\begin{theorem}[vanishing, informal version]\label{thm:van-informal}
For the same character,
\[
                     \eta\circ L_1\Sym^2(p)=0.
\]
\end{theorem}

We construct $\eta$ first, then prove the two theorems in
Sections~\ref{sec:normal} and~\ref{sec:vanishing}.

\section{The quadratic character}

\subsection{The two quadratic complexes}

Let $T=(T_1\xrightarrow{t}T_0)$ be a two-term free resolution of a
finite abelian group $K$.  Thus $t$ is injective and
$\operatorname{coker}t=K$.  The derived symmetric and exterior squares
are computed by
\begin{align}
 K_S(T)&=\left[
 \Lambda^2T_1\xrightarrow{\kappa_2}
 T_1\otimes T_0\xrightarrow{\kappa_1}\Sym^2T_0
 \right],                                                   \label{eq:KS}\\
 K_\Lambda(T)&=\left[
 \Gamma^2T_1\xrightarrow{\lambda_2}
 T_1\otimes T_0\xrightarrow{\lambda_1}\Lambda^2T_0
 \right],                                                   \label{eq:KL}
\end{align}
in homological degrees $2,1,0$.  The differentials used below are
\begin{align*}
 \kappa_2(a\wedge a')&=a\otimes ta'-a'\otimes ta,
 &\kappa_1(a\otimes x)&=ta\,x,\\
 \lambda_2(\gamma_2(a))&=a\otimes ta,
 &\lambda_2(aa')&=a\otimes ta'+a'\otimes ta,\\
 \lambda_1(a\otimes x)&=ta\wedge x.
\end{align*}
Here $\Gamma^2$ denotes the divided square.  Standard Koszul
theory~\cite{Kock} gives
\[
 H_iK_S(T)=L_i\Sym^2(K),
 \qquad H_iK_\Lambda(T)=L_i\Lambda^2(K).
\]

Because $K$ is finite, $t\otimes\Q$ is an isomorphism.  Hence
$K_S(T)\otimes\Q$ is acyclic in positive degrees and all positive
homology of $K_S(T)$ is finite.  The universal coefficient sequence
therefore gives, naturally in $K$,
\begin{equation}\label{eq:UCT}
 H^2\Hom(K_S(T),\Z)
 \cong\Ext^1(L_1\Sym^2(K),\Z)
 \cong\Hom(L_1\Sym^2(K),\Q/\Z).
\end{equation}
Indeed $\Hom(H_2K_S(T),\Z)=0$, and
$\Ext^1(A,\Z)\cong\Hom(A,\Q/\Z)$ for finite $A$.

\subsection{A compatible presentation of the extension}

Choose free resolutions
\[
 P_1\xrightarrow dP_0\longrightarrow V,
 \qquad
 Q_1\xrightarrow eQ_0\longrightarrow W.
\]
Lift $P_0\to V$ to a map $P_0\to H$.  Its composite with $d$ takes
values in $W$; lift that map through $Q_0\to W$ to
$\widetilde b:P_1\to Q_0$.  A compatible free resolution of $H$ is
\begin{equation}\label{eq:R}
 R_1=P_1\oplus Q_1\xrightarrow{\partial}R_0=P_0\oplus Q_0,
 \qquad
 \partial(a,s)=(da,es-\widetilde b(a)).
\end{equation}
Indeed, the map $R_0\to H$ induced by the two chosen lifts has kernel
$\im\partial$, and $\partial$ is injective because $d$ and $e$ are.
Projection $R_\bullet\to P_\bullet$ resolves $p:H\to V$.  If bars
denote passage from $P_0$ to $V$ and from $Q_0$ to $W$, define
\begin{equation}\label{eq:phi}
 \varphi:P_1\otimes P_0\longrightarrow C,
 \qquad
 \varphi(a,x)=g(\bar x,\overline{\widetilde b(a)}).
\end{equation}
Since $\overline{da'}=0$ in $V$,
\[
 \varphi\kappa_2(a\wedge a')
   =\varphi(a,da')-\varphi(a',da)=0.
\]
Thus $\varphi$ is a degree-one cocycle of $\Hom(K_S(P),C)$.

\subsection{The exterior lift}

The following elementary use of bracket antisymmetry is the only point
at which the exterior-square complex enters.

\begin{lemma}[exterior-compatible lift]\label{lem:exterior}
The cocycle $\varphi$ has an integral lift
$\widehat\varphi:P_1\otimes P_0\to\Z$ such that
\begin{equation}\label{eq:exterior}
 \widehat\varphi(a,da')+\widehat\varphi(a',da)=0
 \qquad(a,a'\in P_1).
\end{equation}
Equivalently, $\widehat\varphi$ annihilates
$\im\lambda_2$ in~\eqref{eq:KL}.
\end{lemma}

\begin{proof}
Choose Smith bases
\[
 P_0=\bigoplus_i\Z x_i,
 \qquad P_1=\bigoplus_i\Z a_i,
 \qquad da_i=d_ix_i,
 \qquad d_i\mid d_j\quad(i<j).
\]
Lift the image of $x_i$ to $X_i\in L$.  In $H=L/C$, the element
$\overline{\widetilde b(a_i)}\in W$ is $d_iX_i$ modulo $C$.
Consequently
\begin{equation}\label{eq:phiij}
 \varphi_{ij}:=\varphi(a_i,x_j)
   =[X_j,d_iX_i]=d_i[X_j,X_i]\in C.
\end{equation}
For $i<j$ choose an integer $r_{ij}$ representing $\varphi_{ij}$ under
$C=\frac1q\Z/\Z$; thus
$r_{ij}/q+\Z=\varphi_{ij}$.  Define
\begin{equation}\label{eq:hatphi}
 \widehat\varphi(a_i,x_j)=r_{ij},
 \qquad
 \widehat\varphi(a_j,x_i)=-\frac{d_j}{d_i}r_{ij},
 \qquad
 \widehat\varphi(a_i,x_i)=0.
\end{equation}
The second integer in~\eqref{eq:hatphi} really lifts
$\varphi_{ji}$: by~\eqref{eq:phiij} and antisymmetry,
\[
 \varphi_{ji}=d_j[X_i,X_j]
   =-\frac{d_j}{d_i}\,d_i[X_j,X_i]
   =-\frac{d_j}{d_i}\varphi_{ij}.
\]
Moreover,
\[
 d_j\widehat\varphi(a_i,x_j)
 +d_i\widehat\varphi(a_j,x_i)=0,
 \]
and the diagonal equality is immediate.  Bilinearity gives
\eqref{eq:exterior} for all $a,a'$.
\end{proof}

In particular, the exterior-square Bockstein of $[\varphi]$ is zero:
the chosen integral lift is already a cocycle on $K_\Lambda(P)$.
Since $\varphi(a,da')=0$ in $C$, the integer
$\widehat\varphi(a,da')$ is divisible by $q$.  Hence
\begin{equation}\label{eq:h}
 h(a\wedge a')=\frac{\widehat\varphi(a,da')}{q}
\end{equation}
defines an integral degree-two cochain on $K_S(P)$.  Let
\begin{equation}\label{eq:eta}
 \eta\in\Hom(L_1\Sym^2(V),\Q/\Z)
\end{equation}
($\eta_2$ in the general tower notation) be the character
corresponding through~\eqref{eq:UCT} to the class
$[h]$ (every degree-two cochain is a cocycle because the complex ends
in degree two).  We fix this lift and this character for the rest of
the proof.

For later use, let
$\beta_S:H^1\Hom(K_S(P),C)\to H^2\Hom(K_S(P),\Z)$ be the Bockstein for
$0\to\Z\xrightarrow q\Z\to C\to0$.  It is represented by
\[
 \frac{\widehat\varphi(a,da')-
       \widehat\varphi(a',da)}q
 =2h(a\wedge a')
\]
because of~\eqref{eq:exterior}.  Thus
\begin{equation}\label{eq:twice}
                         \beta_S[\varphi]=2\eta.
\end{equation}
The character $\eta$, rather than the generally insufficient class
$2\eta$, is what the PBW normal form detects.  Thus $\eta$ is the
specific half of the symmetric Bockstein selected by the vanishing of
the exterior Bockstein.

\begin{lemma}[Smith evaluation]\label{lem:evaluation}
For $i<j$ put
\[
 c_{ij}=\frac{d_j}{d_i}a_i\otimes x_j-a_j\otimes x_i
       \in\ker\kappa_1^P.
\]
The $c_{ij}$ form a basis of $\ker\kappa_1^P$, and
$\kappa_2(a_i\wedge a_j)=d_i c_{ij}$.  Therefore, if
$c_u=\sum_{i<j}u_{ij}c_{ij}$, then
\begin{equation}\label{eq:evaluation}
 \eta([c_u])=\frac1q\sum_{i<j}u_{ij}\frac{d_j}{d_i}
                  \widehat\varphi(a_i,x_j)+\Z.
\end{equation}
\end{lemma}

\begin{proof}
The description of $\ker\kappa_1^P$ follows by comparing the
coefficients of $x_i^2$ and $x_ix_j$ in $\Sym^2P_0$.  The displayed
boundary formula is immediate.  Under the universal-coefficient
identification~\eqref{eq:UCT}, the value of the class represented by
$h$ on $[c_{ij}]$ is
\[
 \frac{h(a_i\wedge a_j)}{d_i}+\Z
 =\frac{d_j\widehat\varphi(a_i,x_j)}{qd_i}+\Z.
\]
Linearity gives~\eqref{eq:evaluation}.
\end{proof}

\section{The normal-form theorem}\label{sec:normal}

This section contains the only enveloping-algebra calculation.  We
spell it out because merely writing the four summands of
$K_S(R)_1$ does not prove that an arbitrary dimension witness supplies
all four of them.

Choose bases
\[
 P_0=\bigoplus_i\Z x_i,\quad P_1=\bigoplus_i\Z a_i,
 \quad da_i=d_ix_i,
\]
as in Lemma~\ref{lem:exterior}, and Smith bases
\[
 Q_0=\bigoplus_k\Z y_k,\quad Q_1=\bigoplus_k\Z b_k,
 \quad eb_k=e_ky_k,
\]
with $e_k\mid e_l$ for $k<l$.  Write
\begin{equation}\label{eq:Bmatrix}
       \widetilde b(a_i)=B_i:=\sum_k b_{ik}y_k,
 \qquad D=\operatorname{diag}(d_i),\quad
 E=\operatorname{diag}(e_k),\quad B=(b_{ik}).
\end{equation}

\subsection{The adapted free presentation}

Choose lifts $X_i\in L$ of the images of $x_i$.  They generate $L$:
if $L_0$ is the subring they generate, then $L=L_0+\gm_2(L)$, whence
$\gm_2(L)\subseteq L_0+\gm_3(L)$.  Substituting this last inclusion
once more into $\gm_3(L)=[L,\gm_2(L)]$ gives
$\gm_3(L)\subseteq L_0+\gm_4(L)=L_0$.  Thus $L=L_0$.
Let $F$ be the free Lie ring on symbols $X_i$, mapping onto $L$, and
let $\mathcal R$ be its kernel.  Write
$F=\bigoplus_{r\geq1}F_r$ by bracket length.  Since $\gm_4(L)=0$,
$F_{\geq4}\subseteq\mathcal R$.

Smith reduction of $F_2\to W$ and $F_3\to C$ supplies bases
\[
 F_2=\langle Y_1,\ldots,Y_n,\widetilde Y_1,\ldots\rangle,
 \qquad
 F_3=\langle Z,\widetilde Z_1,\ldots\rangle,
\]
where $Y_k$ maps to the element represented by $y_k$, $Z$ maps to
$z$, and the tilded elements map to zero in $W$ and $C$, respectively.
For suitable integers $c_i,m_k,\widetilde c_l$, the kernel contains
\begin{align}
 R_i&=d_iX_i-\mathcal B_i-c_iZ,
       &\mathcal B_i&=\sum_kb_{ik}Y_k,                    \label{eq:Ri}\\
 S_k&=e_kY_k-m_kZ,
 &\widetilde R_l&=\widetilde Y_l-\widetilde c_lZ,        \label{eq:Sk}\\
 T&=qZ,
 &\widetilde T_l&=\widetilde Z_l.                        \label{eq:Tk}
\end{align}
Together with a basis of $F_4$, these elements generate
$\mathcal R$ modulo $F_{\geq5}$.  Indeed, for a relation, its
degree-one component lies in the kernel of $F_1\to V$ and is removed
by the $R_i$.  Its new degree-two component lies in the kernel of
$F_2\to W$ and is removed by the $S_k,\widetilde R_l$; its new
degree-three component lies in the kernel of $F_3\to C$ and is
removed by $T,\widetilde T_l$.  A basis of $F_4$ removes the last
component.  If the original relation has no component below degree
$r$, the earlier steps are empty, so only relations of least degree at
least $r$ are used.

Write
\begin{equation}\label{eq:Gamma}
 [X_j,Y_k]=g_{jk}Z+\sum_l h_{jkl}\widetilde Z_l,
 \qquad
 \Gamma_{ij}=\sum_kb_{ik}g_{jk}.
\end{equation}
In $C$ this says
\begin{equation}\label{eq:Gamma-phi}
 \Gamma_{ij}/q+\Z=[X_j,d_iX_i]
   =\varphi(a_i,x_j)
   =\widehat\varphi(a_i,x_j)/q+\Z.
\end{equation}

\subsection{PBW extraction}

For a strictly upper-triangular integer matrix $u=(u_{ij})$, set
\begin{equation}\label{eq:ubar}
 \overline u=u-Du^{\mathsf T}D^{-1},\qquad
 \overline u_{ij}=
 \begin{cases}
  u_{ij},&i<j,\\
  -(d_i/d_j)u_{ji},&i>j,\\
  0,&i=j.
 \end{cases}
\end{equation}

\begin{lemma}[PBW coefficient lemma]\label{lem:PBW}
If $\zeta z\in\dlt_5(L)$, there are integer matrices
$u=(u_{ij})_{i<j}$ and $v,v'=(v_{ik}),(v'_{ik})$ such that
\begin{align}
 \overline uB+Dv+v'E&=0,                                 \label{eq:PBW-B}\\
 \zeta&\equiv
  \sum_{i<j}u_{ij}\frac{d_j}{d_i}\Gamma_{ij}\pmod q,    \label{eq:PBW-Z}\\
 (v^{\mathsf T}B+B^{\mathsf T}v)_{kl}
  &\equiv0\pmod{\gcd(e_k,e_l)}\quad(k<l),               \label{eq:PBW-C1}\\
 (v^{\mathsf T}B)_{kk}&\equiv0\pmod{e_k}.               \label{eq:PBW-C2}
\end{align}
\end{lemma}

\begin{proof}
Choose $f\in\gm_3(F)$ mapping to $\zeta z$.  By
\eqref{eq:free-criterion}, there is
$\theta\in\mathfrak r:=U(F)\mathcal R U(F)$ such that
$\theta-f\in\aug(F)^5$.

Give $F_r$ weight $r$ and order homogeneous bases by
\[
 X<Y<\widetilde Y<Z<\widetilde Z<F_4<\cdots.
\]
The PBW theorem implies that the ordered products of total weight
below five form a $\Z$-basis modulo $\aug(F)^5$; here freeness of
$F$ identifies total bracket weight with augmentation degree
\cite[Ch.~I, \S2.7]{Bou}.
The preceding triangular description of $\mathcal R$ shows that,
modulo $\aug(F)^5$, every summand of $\theta$ is a product
$a\sigma b$, where $\sigma$ is one of
$R_i,S_k,\widetilde R_l,T,\widetilde T_l$ or a basis element of
$F_4$, and the possible total weights of $a,b$ are respectively
\[
 0,1,2,3;\qquad 0,1,2;\qquad 0,1;\qquad 0.
\]
This list is exhaustive because these relations have least weights
$1,2,3,4$.

The two identities
\begin{equation}\label{eq:collection-identities}
 h\sigma=\sigma h+[h,\sigma],
 \qquad h'h=hh'-[h,h']\quad(h<h')
\end{equation}
put every such product in PBW order.  The first correction is again a
relation, of at least the sum of the two weights; either correction has
one fewer external factor.  Induction on the number of factors and
then on inversions proves termination without leaving the preceding
list.

Compare the collected expression with the Lie element $f$.
In weight one, independence of the $X_i$ forces every scalar
$R_i$-coefficient to be zero.  In weight two, the diagonal $X_i^2$
coefficient vanishes.  For $i<j$, the only leading packets at
$X_iX_j$ are $R_iX_j$ and $X_iR_j$; if their coefficients are $A,C$,
then $d_iA+d_jC=0$.  Since $d_i\mid d_j$, all the remaining
weight-two packets are therefore
\begin{equation}\label{eq:Theta2}
 \sum_{i<j}u_{ij}
 \left(\frac{d_j}{d_i}R_iX_j-X_iR_j\right).
\end{equation}
Scalar $S_k$- and $\widetilde R_l$-terms also vanish, because their
leading $e_kY_k$ and $\widetilde Y_l$ terms are independent.

At weight three, retain from the collected expression the packets
$R_iY_k$ and $X_iS_k$, with coefficients $v_{ik}$ and $v'_{ik}$.
The other canonical packets are
$R_i\widetilde Y_l$, $X_i\widetilde R_l$, $T$,
$\widetilde T_l$, and the three ordered placements of one $R_i$ with
two $X$-factors.  Their leading terms show that they have no named
$X_iY_k$ coefficient; only $T=qZ$ can have a named one-factor
$Z$ coefficient at this weight.
Put
\begin{equation}\label{eq:Theta}
 \Theta=\sum_{i<j}u_{ij}
   \left(\frac{d_j}{d_i}R_iX_j-X_iR_j\right)
   +\sum_{i,k}v_{ik}R_iY_k
   +\sum_{i,k}v'_{ik}X_iS_k.
\end{equation}
Expanding
\eqref{eq:Ri}--\eqref{eq:Sk} and using
$Y_kX_j=X_jY_k-[X_j,Y_k]$ gives
\begin{align}
 [X_iY_k]\Theta&=(\overline uB+Dv+v'E)_{ik},              \label{eq:coeff-XY}\\
 [Z]\Theta&=\sum_{i<j}u_{ij}\frac{d_j}{d_i}\Gamma_{ij}, \label{eq:coeff-Z}\\
 [Y_kY_l]\Theta
   &=-(v^{\mathsf T}B+B^{\mathsf T}v)_{kl} &&(k<l),      \label{eq:coeff-YY}\\
 [Y_k^2]\Theta&=-(v^{\mathsf T}B)_{kk}.                  \label{eq:coeff-Y2}
\end{align}
Here $[M]\Theta$ denotes the coefficient of the PBW monomial $M$.
For example, the $i,j$ packet in~\eqref{eq:Theta2} contains
$-(d_j/d_i)\mathcal B_iX_j+X_i\mathcal B_j$; its first term contributes
$-(d_j/d_i)b_{ik}$ at $X_jY_k$ and
$(d_j/d_i)\Gamma_{ij}$ at $Z$.  This verifies both the sign and the
indices in~\eqref{eq:coeff-XY}--\eqref{eq:coeff-Z}.  The last two
lines come from $-\sum v_{ik}\mathcal B_iY_k$.

For completeness, all other packets of total weight four have the
following leading shapes.  A letter $Y$ or $Z$ may be named or
supplementary.
\[
\begin{array}{c|c|c|c}
 \text{relation}&\text{external factors}&\text{leading shape}
      &\text{named }Y_kY_l\\ \hline
 R_i&XXX&XXXX&0\\
 R_i&X,Y&XXY&0\\
 R_i&Z&XZ&0\\
 S_k,\widetilde R_l&XX&XXY&0\\
 \widetilde R_l&Y&\widetilde Y_lY&0\\
 T,\widetilde T_l&X&XZ&0\\
 F_4&1&F_4&0.
\end{array}
\]
The omitted row is $S_k$ with one named $Y$-factor:
for $k<l$ its two placements contribute respectively
$e_kY_kY_l$ and $e_lY_kY_l$, and on the diagonal they contribute
$e_kY_k^2$.  The table is closed under both operations in
\eqref{eq:collection-identities}: a correction either replaces two
$X$'s by one $Y$, replaces $X,Y$ by one $Z$, or raises the least
weight of the relation.  Thus it also accounts for every correction
created during collection.

No omitted packet has an $X_iY_k$ coefficient.  Its $Z$ coefficient
is a multiple of $q$ (or is supplementary), and its named
$Y_kY_l$ coefficient lies in
$e_k\Z+e_l\Z=\gcd(e_k,e_l)\Z$; on the diagonal it lies in
$e_k\Z$.  The Lie element $f$ has no multi-factor PBW coefficient,
and its $Z$-coefficient is $\zeta$ modulo $q$.  Comparing
\eqref{eq:coeff-XY}--\eqref{eq:coeff-Y2} with
$\theta-f\in\aug(F)^5$ proves
\eqref{eq:PBW-B}--\eqref{eq:PBW-C2}.
\end{proof}

We now translate the coefficient lemma into the promised functorial
normal form.  Decompose
\[
 K_S(R)_1=(P_1\otimes P_0)\oplus(P_1\otimes Q_0)
 \oplus(Q_1\otimes P_0)\oplus(Q_1\otimes Q_0),
\]
let $\tau$ interchange factors, and let $\mu$ be symmetric
multiplication in $\Sym^2Q_0$.

\begin{lemma}[cycle extracted from PBW]\label{lem:PBW-cycle}
For the matrices supplied by Lemma~\ref{lem:PBW}, set
\[
 c=c_u,\qquad
 \nu=\sum_{i,k}v_{ik}a_i\otimes y_k,\qquad
 \nu'=\sum_{i,k}v'_{ik}b_k\otimes x_i.
\]
There is $w\in Q_1\otimes Q_0$ such that
$\chi=c+\nu+\nu'+w$ is a cycle in $K_S(R)$, and
\begin{equation}\label{eq:Zledger}
 \zeta\equiv
 \sum_{i<j}u_{ij}\frac{d_j}{d_i}
       \widehat\varphi(a_i,x_j)\pmod q.
\end{equation}
\end{lemma}

\begin{proof}
Equation~\eqref{eq:PBW-B} is exactly
\[
 -\tau(\widetilde b\otimes1)c+(d\otimes1)\nu
       +\tau(e\otimes1)\nu'=0.
\]
Equations~\eqref{eq:PBW-C1}--\eqref{eq:PBW-C2} say exactly that
$\mu(\widetilde b\otimes1)\nu$ is zero in
\[
 \Sym^2(W)=\operatorname{coker}
   (\kappa_1^Q:Q_1\otimes Q_0\to\Sym^2Q_0),
\]
because the orders of $y_ky_l$ and $y_k^2$ there are
$\gcd(e_k,e_l)$ and $e_k$.  Hence $w$ may be chosen so that
\[
 -\mu(\widetilde b\otimes1)\nu+\kappa_1^Q(w)=0.
\]
Direct substitution in~\eqref{eq:R} gives
\[
 \kappa_1^R(\chi)=
 \kappa_1^P(c)\ \oplus\
 \bigl[-\tau(\widetilde b\otimes1)c+(d\otimes1)\nu
          +\tau(e\otimes1)\nu'\bigr]\ \oplus\
 \bigl[-\mu(\widetilde b\otimes1)\nu+\kappa_1^Q(w)\bigr].
\]
All three components vanish, so $\chi$ is a cycle.
Finally,~\eqref{eq:Zledger} follows from
\eqref{eq:PBW-Z} and~\eqref{eq:Gamma-phi}.
\end{proof}

\begin{theorem}[normal form]\label{thm:normal}
If $\zeta z\in\dlt_5(L)$, there is
$\xi\in L_1\Sym^2(H)$ such that
\begin{equation}\label{eq:normal}
 \frac{\zeta}{q}+\Z
 =\eta\bigl(L_1\Sym^2(p)(\xi)\bigr).
\end{equation}
\end{theorem}

\begin{proof}
Take the cycle $\chi$ from Lemma~\ref{lem:PBW-cycle} and put
$\xi=[\chi]\in H_1K_S(R)=L_1\Sym^2(H)$.  Projection
$R_\bullet\to P_\bullet$ sends $\chi$ to $c_u$, hence
\[
 L_1\Sym^2(p)(\xi)=[c_u].
\]
Now divide~\eqref{eq:Zledger} by $q$ and apply
Lemma~\ref{lem:evaluation}.
\end{proof}

The normal form is a necessary consequence of membership in
$\dlt_5(L)$; no converse is asserted.

\section{The vanishing theorem}\label{sec:vanishing}

The proof in this section uses only the additive extension
\eqref{eq:extension}, the pairing~\eqref{eq:g}, and the exterior lift
of Lemma~\ref{lem:exterior}.  It is independent of PBW collection.

Put
\[
 A=\Hom(W,C),
 \qquad A^\vee=\Hom(A,C).
\]
Evaluation gives a natural isomorphism
\begin{equation}\label{eq:bidual}
                         W/qW\xrightarrow{\sim}A^\vee.
\end{equation}
For $W=\Z/e$, the class of $1$ evaluates the generator of
$\Hom(W,C)$ onto an element of order $\gcd(e,q)$; thus the evaluation
map is an isomorphism between two cyclic groups of that order.  The
general case follows by decomposing $W$ into cyclic summands.
Define
\begin{equation}\label{eq:rho-b}
 \rho:P_0\to A,
 \quad \rho(x)(w)=g(\bar x,w),
 \qquad
 b:P_1\to A^\vee,
 \quad b(a)=\overline{\widetilde b(a)}\pmod{qW}.
\end{equation}
Then $\rho d=0$ and
\begin{equation}\label{eq:factorphi}
                         \varphi(a,x)=b(a)(\rho x).
\end{equation}

\begin{lemma}[polarisation]\label{lem:polarisation}
If $G_0$ is a finite abelian group and
$\Omega:G_0\times G_0\to\Q/\Z$ is alternating, there is a bilinear
$H:G_0\times G_0\to\Q/\Z$ such that $H^{\mathsf T}-H=\Omega$.
\end{lemma}

\begin{proof}
Choose cyclic generators $e_1,\ldots,e_r$ of $G_0$.  For $i<j$, set
$H(e_j,e_i)=\Omega(e_i,e_j)$ and $H(e_i,e_j)=H(e_i,e_i)=0$.
The orders of both generators annihilate each assigned value, so $H$
is well defined and bilinear.  On every pair of generators,
$H^{\mathsf T}-H=\Omega$.
\end{proof}

\begin{proposition}[quadratic certificate]\label{prop:certificate}
There exist a homomorphism $\alpha:P_0\to A$, a symmetric bilinear form
\[
 \lambda:A^\vee\times A^\vee\longrightarrow C,
\]
and a cocycle
\[
 \psi(a,x)=b(a)(\alpha x)
     \in\Hom(P_1\otimes P_0,C)
\]
such that
\begin{equation}\label{eq:certificate}
              \alpha d=\lambda_*b,
 \qquad
              \beta_S[\psi]=\eta.
\end{equation}
Here $(\lambda_*\chi)(\chi')=\lambda(\chi,\chi')$ and
$A\cong\Hom(A^\vee,C)$ by evaluation.
\end{proposition}

\begin{proof}
Form the pushout
\begin{equation}\label{eq:E}
 \mathcal E=
 \frac{A^\vee\oplus P_0}
      {\langle(b(a),-da):a\in P_1\rangle}.
\end{equation}
Since $d$ is injective, the evident maps give an exact sequence
$0\to A^\vee\to\mathcal E\to V\to0$: if $[\chi,0]=0$, then
$(\chi,0)=(b(a),-da)$ for some $a$, hence $a=0$ and $\chi=0$; after
quotienting by $A^\vee$, the cokernel is $P_0/dP_1=V$.  Write
$i(\chi)=[\chi,0]$ and $j(x)=[0,x]$; the defining relation says
\begin{equation}\label{eq:jdib}
                         j(da)=i(b(a)).
\end{equation}

Because $d\otimes\Q$ is an isomorphism, the rule
\begin{equation}\label{eq:vartheta}
 \vartheta(da,x)=\frac{\widehat\varphi(a,x)}q
\end{equation}
defines a rational bilinear form on $P_0\otimes\Q$.
Equation~\eqref{eq:exterior} says that $\vartheta$ is alternating.
Define
\begin{equation}\label{eq:Omega}
 \Omega([\chi,x],[\chi',x'])
 =\chi(\rho x')-\chi'(\rho x)+\vartheta(x,x')+\Z.
\end{equation}
This is well defined on $\mathcal E$.  Indeed, in the first variable,
\begin{align*}
 \Omega((b(a),-da),(\chi',x'))
 &=b(a)(\rho x')-\vartheta(da,x')+\Z\\
 &=\varphi(a,x')-\frac{\widehat\varphi(a,x')}q+\Z=0,
\end{align*}
and alternation gives the same conclusion in the second variable.
The form $\Omega$ is alternating.

Choose $H$ as in Lemma~\ref{lem:polarisation} and define
\begin{equation}\label{eq:alpha-lambda}
 \chi(\alpha x)=H(jx,i\chi),
 \qquad
 \lambda(\chi,\chi')=H(i\chi',i\chi).
\end{equation}
These values lie in $C$, since $i(A^\vee)$ is killed by $q$.  The
restriction of $\Omega$ to $i(A^\vee)^2$ is zero, so $H$ is symmetric
there and $\lambda$ is symmetric.  By~\eqref{eq:jdib},
\[
 \chi(\alpha da)=H(i b(a),i\chi)
 =H(i\chi,i b(a))=\lambda(b(a),\chi),
\]
which proves $\alpha d=\lambda_*b$.  Symmetry of $\lambda$ also shows
directly that $\psi$ is a cocycle:
\[
 \psi(a,da')-\psi(a',da)
 =\lambda(b(a'),b(a))-\lambda(b(a),b(a'))=0.
\]

Choose a rational bilinear lift $P$ of
$(x,x')\mapsto H(jx,jx')$.  Initially
$P^{\mathsf T}-P-\vartheta$ is integral and alternating.  Adding an
integral bilinear form supported in one triangle, we may and do arrange
\begin{equation}\label{eq:P}
                         P^{\mathsf T}-P=\vartheta.
\end{equation}
Set
\[
                  \widehat\psi(a,x)=qP(x,da).
\]
This is integer-valued: modulo $\Z$,
\[
 P(x,da)=H(jx,jda)=H(jx,ib(a))=b(a)(\alpha x)=\psi(a,x),
\]
and $\psi(a,x)$ lies in $C=\frac1q\Z/\Z$.  The same equality shows
that $\widehat\psi$ is an integral lift of $\psi$.  Finally,
\begin{align*}
 \widehat\psi(a,da')-\widehat\psi(a',da)
 &=q\bigl(P(da',da)-P(da,da')\bigr)\\
 &=q\vartheta(da,da')
 =\widehat\varphi(a,da').
\end{align*}
After division by $q$, the Bockstein cocycle for $\psi$ is exactly the
cochain $h$ of~\eqref{eq:h}.  Hence $\beta_S[\psi]=\eta$.
\end{proof}

\begin{theorem}[vanishing]\label{thm:vanishing}
For the character~\eqref{eq:eta},
\begin{equation}\label{eq:vanishing}
                       \eta\circ L_1\Sym^2(p)=0.
\end{equation}
\end{theorem}

\begin{proof}
Let $f:K_S(R)\to K_S(P)$ be induced by the projection
$R_\bullet\to P_\bullet$.  If $y\in Q_0$, write $y_W$ for its image in
$W$ and $\bar y$ for its image in $W/qW=A^\vee$.  Define a symmetric
$C$-valued form on $R_0=P_0\oplus Q_0$ by
\begin{equation}\label{eq:T}
\begin{split}
 T((x,y),(x',y'))={}&-\alpha(x)(y'_W)-\alpha(x')(y_W)\\
                    &-\lambda(\bar y,\bar y').
\end{split}
\end{equation}
For $(a,s)\in R_1$, equations~\eqref{eq:R} and
\eqref{eq:certificate} give
\begin{align*}
 T(\partial(a,s),(x,y))
 &=-\alpha(da)(y_W)
   +\alpha(x)(\overline{\widetilde b(a)})
   +\lambda(b(a),\bar y)\\
 &=-\lambda(b(a),\bar y)+\psi(a,x)
   +\lambda(b(a),\bar y)\\
 &=\psi(a,x).
\end{align*}
Thus the pulled-back cocycle $f^*\psi$ is the coboundary of $T$, so
$f^*[\psi]=0$.  Naturality of the Bockstein and
Proposition~\ref{prop:certificate} now give
\[
 f^*\eta=f^*\beta_S[\psi]
 =\beta_S(f^*[\psi])=0.
\]
Under the natural identification~\eqref{eq:UCT}, $f^*$ is
precomposition with $L_1\Sym^2(p)$.  This is exactly
\eqref{eq:vanishing}.
\end{proof}

\section{Conclusion and the unreduced statement}

\begin{proof}[Proof of Theorem~\ref{thm:reduced}]
Let $a\in\dlt_5(L)$.  Since the published equality
$\dlt_3(L)=\gm_3(L)$ gives
$\dlt_5(L)\subseteq\dlt_3(L)=C$, write $a=\zeta z$.
The normal-form theorem supplies $\xi\in L_1\Sym^2(H)$ with
\[
 \frac\zeta q+\Z
 =\eta(L_1\Sym^2(p)(\xi)).
\]
The right side is zero by Theorem~\ref{thm:vanishing}.  Hence
$q\mid\zeta$, and therefore $a=\zeta z=0$ in $C$.
\end{proof}

We finish by recording why this reduced theorem gives the usual
statement for arbitrary Lie rings.

\begin{corollary}\label{cor:global}
For every Lie ring $L$,
\[
                         \dlt_5(L)\subseteq\gm_4(L).
\]
\end{corollary}

\begin{proof}
Pass to $\bar L=L/\gm_4(L)$.  It is enough, by functoriality of
dimension subrings, to prove $\dlt_5(\bar L)=0$.  Suppose that
$0\ne a\in\dlt_5(\bar L)$.  The ring $\bar L$ has class at most three
and is therefore metabelian.  The published metabelian exponent
theorem~\cite[Theorem~2.1]{PS} asserts
$2\dlt_r(K)\subseteq\gm_r(K)$ for every metabelian $K$; applied to
$\bar L$, it gives $2a=0$, because
$\gm_5(\bar L)=0$.

The enveloping-algebra functor, being a left adjoint, commutes with
filtered colimits.  The equality expressing $a$ as an element of the
fifth augmentation power is finite, so it already holds at some
finitely generated Lie subring $M\subseteq\bar L$.  Thus
$0\ne a\in\dlt_5(M)$.  Since $M$ is finitely generated and has class
at most three, finitely many basic commutators of weights at most three
generate its additive group.  The structure theorem for finitely
generated abelian groups shows that, for all sufficiently large $N$,
the nonzero element $a$ of order two survives in $M/2^NM$; this is a finite
Lie ring whose additive group is a $2$-group.

In this quotient, the image of $a$ lies in
$\dlt_5\subseteq\dlt_3=\gm_3$.  Its third lower-central term is a
finite central abelian $2$-group.
Decompose that group into cyclic summands and project onto a summand on
which the image of $a$ is nonzero.  The kernel is a central ideal.
By functoriality, after quotienting by it, $a$ is still a nonzero
element of $\dlt_5$,
$\gm_4=0$, and $\gm_3$ is cyclic.  This contradicts
Theorem~\ref{thm:reduced}.
\end{proof}

\begin{remark}
The finiteness and $2$-primary hypotheses in
Theorem~\ref{thm:reduced} are therefore not extra assumptions in the
global theorem: they are consequences of the preceding counterexample
reduction.  They are stated in the reduced theorem because the Smith
resolutions, Pontryagin duality~\eqref{eq:bidual}, and the polarisation
argument use them explicitly.
\end{remark}

\begin{thebibliography}{9}

\bibitem{BP}
L.~Bartholdi and I.~B.~S.~Passi,
\emph{Lie dimension subrings},
Internat. J. Algebra Comput. \textbf{25} (2015), 1301--1325.

\bibitem{Bou}
N.~Bourbaki,
\emph{Lie Groups and Lie Algebras, Chapters 1--3},
Springer, 1989.

\bibitem{Kock}
B.~K\"ock,
\emph{Computing the homology of Koszul complexes},
Trans. Amer. Math. Soc. \textbf{353} (2001), 3115--3147.

\bibitem{PS}
I.~B.~S.~Passi and T.~Sicking,
\emph{Dimension quotients of metabelian Lie rings},
Internat. J. Algebra Comput. \textbf{27} (2017), 251--258.

\end{thebibliography}

\end{document}
